\Capitulo{Proyección metodológica}

\subsection{Perspectiva epistemológica}\label{sec:perspectiva-epistemologica}

La perspectiva epistemológica del presente estudio se sustenta en la noción del paradigma científico formulado por \textcite{Kuhn1971revoluciones}, quien entiende por paradigmas aquellas realizaciones científicas universalmente reconocidas que, durante cierto tiempo, proporcionan modelos de problemas y soluciones a una comunidad. Para el autor, un paradigma no se reduce a un cuerpo de teorías, comprende también los métodos, los instrumentos, los criterios de validez y el conjunto de problemas que la comunidad científica reconoce como legítimos en un momento dado.

Desde esta postura se realiza un análisis sistémico del objeto de estudio, que examina las Herramientas de Computación Generativa como totalidad y, dentro de ella, los procedimientos y técnicas para diseñar y validar arquitecturas de autorreflexión sobre \textit{pipelines} RAG como las partes que la integran. El análisis del todo permite situar el objeto en el estado actual de la disciplina; el de las partes, esto es, el mecanismo de recuperación, el componente crítico y el generador, permite localizar con precisión el vacío que la investigación aborda, delimitado al desempeño en dominio específico empresarial sobre modelos de lenguaje pequeños.

Dentro de ese paradigma, la presente investigación se sitúa en lo que Kuhn denomina ciencia normal, no pretende sustituir la matriz disciplinar vigente, sino resolver una anomalía que esta no explica satisfactoriamente en su formulación actual. Dicha anomalía se manifiesta en la generalización deficiente del componente crítico cuando se traslada a un dominio específico (\cite{jeong_improving_2024}), y en que la solución disponible ---reentrenar el crítico con anotaciones obtenidas de un modelo propietario--- resulta inaplicable bajo las restricciones de privacidad de la información e infraestructura de cómputo propias de una pequeña o mediana empresa. A partir del resultado de este análisis se elabora la propuesta de solución mediante el Modelo crítico híbrido, Critic-EmpRAG1.0.

En cuanto a la relación sujeto-objeto, el investigador interactúa con el objeto de estudio sin formar parte de él, interviene deliberadamente mediante la introducción del modelo crítico híbrido y observa el efecto que dicha intervención produce. Esta interacción da paso a una validación empírica organizada en dos componentes: un proceso preexperimental que contrasta los contextos $A_0$ y $A_1$, y una comparación controlada entre configuraciones del sistema. En ambos, el tratamiento de los datos se realiza con herramientas cuantitativas de la estadística inferencial, conforme se detalla en el epígrafe de análisis de los datos. La validación teórica de la propuesta que se realiza mediante la consulta a expertos aplicando el método Delphi. Ambos componentes resultan complementarios y consistentes con el enfoque mixto de la investigación.

\subsection{Enfoque, diseño y tipo de investigación}\label{sec:enfoque-diseno-tipo}

El tipo de investigación es propositiva y aplicada. A partir del análisis de la situación actual en el contexto donde se investiga, se propone como solución, una arquitectura de autorreflexión integrada a un pipeline RAG sobre modelos de lenguaje pequeños; como elemento diferenciador la propuesta del Modelo crítico híbrido, Critic-EmpRAG1.0, alineado a mejorar el acceso al Sistema de soporte de conocimiento empresarial interno en las áreas de TI y desarrollo de Dualbiz SRL. Una vez diseñada e implementada la arquitectura, esta se llevará a la práctica mediante un prototipo para verificar la transformación en la productividad y el acceso al Sistema de soporte de conocimiento empresarial interno como meta de la investigación. El alcance es explicativo, en tanto se manipula deliberadamente la variable independiente para establecer su efecto sobre la variable dependiente. La validación empírica se organiza en dos niveles: un preexperimento de preprueba y posprueba con un solo grupo, que contrasta el acceso al Sistema de conocimiento en el contexto inicial $A_0$ y final $A_1$ después de la aplicación del Modelo crítico híbrido, Critic-EmpRAG1.0; y una comparación controlada entre configuraciones del sistema, que permite atribuir el efecto observado al componente que lo produce. La temporalidad del estudio es longitudinal prospectiva en el tiempo, dado que la validación requiere observar el desempeño del prototipo a lo largo de las iteraciones del proceso experimental.

El enfoque de la presente investigación es mixto. La pregunta de investigación requiere, mejorar la productividad y el acceso al Sistema de soporte de conocimiento empresarial interno para las áreas de TI y desarrollo en Dualbiz SRL, mediante Arquitecturas de autorreflexión y RAG, lo que exige un componente cuantitativo. Por otra parte, se profundiza en las percepciones y necesidades del personal de Dualbiz SRL respecto al acceso al conocimiento operativo interno, lo que exige un componente cualitativo.

\subsection{Delimitación de la investigación}

\subsubsection*{Delimitación temática}

El trabajo de investigación se delimita al diseño, implementación y evaluación de arquitecturas de autorreflexión integradas a pipelines de RAG, integradas en un Modelo (las variables/categorías ya enmarcadas en el Objeto de estudio y el Campo de acción). Se toma como marco teórico y referencial de partida, los aportes de \textcite{Lewis2020rag} para RAG, \textcite{Shinn2023_Reflexion} para autorreflexión mediante refuerzo verbal y \textcite{Asai2024_Selfrag} para la integración nativa de ambos mecanismos (Self-RAG), entre otros autores identificados en el Estudio de Mapeo Sistemático conducido como fase previa de este trabajo (\cite{Kitchenham2007guidelines}; \cite{Petersen2008systematic}).

El estudio se circunscribe al área de conocimiento de la Inteligencia Artificial aplicada al Procesamiento de Lenguaje Natural, específicamente a modelos de lenguaje pequeños (SLM, en el rango de 1B a 13B parámetros). Quedan fuera del alcance el preentrenamiento o ajuste fino de nuevos modelos base, así como los enfoques de Computación Generativa con base en GAN, explorados en el SMS. Esta exclusión obedece a que la propuesta se orienta a la composición de mecanismos ya validados en la literatura (RAG y autorreflexión) sobre modelos de código abierto existentes, y traslada la especificidad de dominio al conocimiento recuperado en tiempo de inferencia, sin recurrir a anotaciones de un modelo propietario ni al reentrenamiento sobre datos de la organización. En cuanto a normas o regulaciones, el diseño del prototipo y el tratamiento del conocimiento operativo de Dualbiz SRL se enmarcan en las políticas internas de confidencialidad de la información, al no requerir el envío de información privada a servicios de terceros.

El despliegue local elimina la exposición a servicios de terceros, pero no agota el tratamiento de la confidencialidad: conserva superficies de exposición propias que el estudio identifica y somete a control. Estas comprenden el índice vectorial y las representaciones densas derivadas del corpus, de las cuales es posible reconstruir parcialmente el texto de origen; los registros de consulta generados en operación; los respaldos del índice y del corpus documental; y las credenciales y permisos de acceso diferenciados por rol. Para cada una se define el control aplicable ---cifrado en reposo, segmentación de accesos, política de retención y anonimización del banco de evaluación---, y la proporción de superficies con control implementado se registra como indicador conforme a la Tabla \ref{tab:operacionalizacion}.

\subsubsection*{Delimitación espacial}

El presente trabajo se lleva a cabo, para su validación empírica, en la empresa Dualbiz SRL, localizada en Santa Cruz de la Sierra, Bolivia, afectando el proceso de soporte y acceso al conocimiento operativo del personal (nuevo y existente, egresados y/o profesionales informáticos y de sistemas) de las áreas de desarrollo y soporte. Si bien el estudio se realiza y valida sobre el conocimiento operativo específico de dicha organización, los resultados y la arquitectura que se propone podrían transferirse a otras pequeñas y medianas empresas de desarrollo de \textit{software} que compartan las condiciones de contorno del caso estudiado ---volumen y naturaleza del corpus documental, perfil del personal consultante y restricciones de privacidad de datos e infraestructura de cómputo---, transferencia que se plantea como hipótesis sujeta a validación posterior y no como resultado comprometido por el presente estudio.

\subsubsection*{Delimitación temporal}

La investigación se desarrolla en cuatro hitos consecutivos, alineados a los objetivos específicos: (1) Sistematizar la literatura acerca de los métodos, técnicas y despliegues de Arquitecturas de autorreflexión y generación aumentada por recuperación (RAG), (2) Diagnosticar el estado actual de los procesos en el acceso al Sistema de soporte de conocimiento empresarial interno en las áreas de TI y desarrollo en Dualbiz SRL, (3) Desarrollar el Modelo crítico híbrido, Critic-EmpRAG1.0, mediante un prototipo funcional sobre modelos de lenguaje pequeños, (4) Validar teórica y empíricamente el Modelo crítico híbrido, Critic-EmpRAG1.0, en sus niveles organizacional y computacional. Se plantea como inicio de cronograma de actividades en octubre del 2026 y fecha de finalización octubre del 2027; este mismo se detalla con sus hitos para cada etapa en la Tabla \ref{tab:cronograma}.

Un horizonte de 12 meses no permite observar directamente la evolución del modelo ni su mantenimiento a escala multianual. Para no renunciar por ello a la evidencia sobre su sostenibilidad en el tiempo, el estudio delimita dicha propiedad a lo que resulta verificable dentro de la ventana temporal disponible. Al residir el conocimiento del dominio en el índice documental y no en los pesos del modelo, y al no requerir el modelo crítico híbrido reentrenamiento alguno sobre datos propietarios, el mantenimiento del sistema se reduce a operaciones acotadas de indexación y sustitución de componentes, y no a un ciclo de reentrenamiento. Por lo que, la escalabilidad y la evolución del modelo se someten a las siguientes pruebas, incorporadas a la etapa de validación: Escalabilidad respecto al volumen de conocimiento, Evolución y vigencia del conocimiento, Mantenimiento evolutivo del modelo base, entre otros aspectos.

Quedan fuera del alcance temporal de esta investigación la observación de la adopción organizacional sostenida del prototipo, la estimación del retorno de la inversión en nuevos equipos u otra infraestructura fuera de la planteada en el presente contexto, o el cambio a otros dominios y el mantenimiento correctivo en operación productiva posterior a la defensa. Dichos aspectos se enuncian como líneas de continuidad y no como resultados comprometidos por el presente estudio.

\subsection{Variables de Investigación}

\subsubsection{Definición Conceptual}

\subsubsection*{Variable independiente}

\textbf{Modelo crítico híbrido Critic-EmpRAG1.0.} Constructo que propone la presente investigación. Se define como el componente de una arquitectura RAG que evalúa la evidencia recuperada del corpus de conocimiento, juzga su correspondencia con la consulta y con la respuesta en elaboración, y desencadena ciclos correctivos de recuperación o regeneración antes de emitir la respuesta. Recoge de la literatura sobre autorreflexión la noción de un juicio que retroalimenta al propio sistema (\cite{Shinn2023_Reflexion}) y, de Self-RAG, la de un crítico integrado al proceso de recuperación y generación (\cite{Asai2024_Selfrag}). Se distingue de ambos en la fuente de su conocimiento del dominio: en lugar de residir en parámetros ajustados con anotaciones de un modelo propietario (\cite{jeong_improving_2024}), reside en la evidencia recuperada en tiempo de inferencia, con lo que extiende al componente crítico la distinción entre conocimiento paramétrico y no paramétrico que \textcite{Lewis2020rag} formularon para el generador. En este caso particular, se aplica a bases de conocimiento empresarial de la industria del \textit{software}.

\subsubsection*{Variable dependiente 1}

\textbf{Productividad del proceso de consulta.} \textcite{Drucker1999_KWProductivity} sostiene que la productividad del trabajador del conocimiento no se reduce a la cantidad de resultados por unidad de tiempo: exige considerar la calidad del resultado y depende de la autonomía con que el trabajador gestiona su propia tarea. En coherencia con ese planteamiento, se define como el grado en que el personal de las áreas de TI y desarrollo resuelve sus consultas operativas en menos tiempo, con menor dependencia del personal experimentado y con menor esfuerzo. Cuenta como referente empírico la evidencia de campo de \textcite{Brynjolfsson2025_GenAIWork}, quienes miden la productividad de un proceso de soporte como consultas resueltas por hora y reportan que el efecto de los asistentes generativos se concentra en el personal con menor experiencia. Se descompone en tres dimensiones:

\begin{itemize}
	\item \textbf{Eficiencia.} Tiempo necesario para resolver una consulta operativa.
	\item \textbf{Autonomía del personal.} Capacidad del personal para resolver sus consultas sin recurrir a la supervisión experimentada, en el sentido que \textcite{Drucker1999_KWProductivity} atribuye a la autogestión del trabajador del conocimiento, y con incidencia directa sobre los procesos de inducción e incorporación del personal (\cite{Achury2016induccion}).
	\item \textbf{Esfuerzo.} Carga que la consulta impone al consultante, según su propia percepción.
\end{itemize}

\subsubsection*{Variable dependiente 2}

\textbf{Acceso al conocimiento.} \textcite{Alavi2001_KMS} sitúan el almacenamiento, la recuperación y la transferencia del conocimiento entre los procesos centrales que los sistemas de gestión del conocimiento deben soportar en una organización. Sobre esa base, se define como el grado en que el personal de las áreas de TI y desarrollo obtiene del Sistema de soporte de conocimiento empresarial interno respuestas correctas, fundamentadas en evidencia verificable y entregadas dentro de las restricciones de privacidad de la información e infraestructura de cómputo de la organización. Se descompone en dos dimensiones:

\begin{itemize}
	\item \textbf{Fundamentación y confiabilidad.} Grado en que las respuestas se apoyan en evidencia recuperada verificable y se encuentran libres de afirmaciones no sustentadas. Es la dimensión en la que opera el mecanismo postulado en la hipótesis y se sustenta en la noción de modelos de lenguaje fiables y atribuibles de \textcite{Asai2024_RALMs}, así como en la evidencia de que el anclaje en evidencia recuperada reduce la alucinación (\cite{Shuster2021_Hallucination}).
	\item \textbf{Viabilidad bajo restricciones.} Grado en que el sistema entrega las respuestas dentro de los límites de infraestructura de cómputo disponibles y sin exponer conocimiento confidencial a servicios de terceros, condición que el Estudio de Mapeo Sistemático conducido como fase previa identifica como vacío persistente del campo (\cite{Alcides2026}). Se incluye en el acceso al conocimiento porque una respuesta que no llega a tiempo o que exige exponer información confidencial no constituye un acceso utilizable en la organización.
\end{itemize}


\subsubsection{Definición operacional de las variables}

La operacionalización de las variables se desglosa en dimensiones e indicadores según las definiciones conceptuales anteriores, conforme se presenta en la Tabla \ref{tab:operacionalizacion}.

\begin{table}[H]
\caption{Operacionalización de las variables}
\label{tab:operacionalizacion}
\setstretch{1}
\renewcommand{\arraystretch}{1.3}
\small
\begin{tabularx}{\textwidth}{|L{3.2cm}|L{3.4cm}|Y|}
\hline
\rowcolor{gray!25}
\textbf{Variable} & \textbf{Dimensiones} & \textbf{Indicadores} \\
\hline
1. VI. Modelo crítico Critic-EmpRAG1.0
 & 1.1. Mecanismo de recuperación (RAG)
 & 1.1.1. Estrategia de recuperación (densa, híbrida o adaptativa por incertidumbre) \newline
   1.1.2. Condición de activación de la recuperación (siempre activa vs.\ condicional) \\
\cline{2-3}
 & 1.2. Mecanismo de autorreflexión
 & 1.2.1. Tipo de autocrítica (verbal, \textit{tokens} aprendidos o basada en reglas) \newline
   1.2.2. Número máximo de ciclos de refinamiento iterativo \newline
   1.2.3. Fuente de verificación de la crítica (evidencia externa recuperada vs.\ conocimiento paramétrico) \\
\hline
2. VD. Productividad y acceso al conocimiento
 & 2.1. Fundamentación y confiabilidad
 & 2.1.1. Exactitud de respuesta (\textit{accuracy}/F1 frente a la respuesta esperada) \newline
   2.1.2. Tasa de alucinación (afirmaciones no sustentadas por la evidencia recuperada, sobre el total de afirmaciones verificables) \newline
   2.1.3. Proporción de afirmaciones con fuente o cita verificable \newline
   2.1.4. Puntaje de fidelidad (\textit{faithfulness}) respecto a los documentos recuperados \newline
   2.1.5. Tasa de abstención correcta ante consultas sin respaldo en el corpus \\
\cline{2-3}
 & 2.2. Productividad del proceso de consulta
 & 2.2.1. Tiempo medio de resolución por consulta operativa \newline
   2.2.2. Tasa de consultas resueltas sin escalamiento a personal experimentado \newline
   2.2.3. Interrupciones semanales a personal experimentado por consultas de conocimiento \newline
   2.2.4. Esfuerzo percibido en la consulta (escala Likert de 1 a 5) \\
\cline{2-3}
 & 2.3. Viabilidad bajo restricciones
 & 2.3.1. Latencia promedio de respuesta \newline
   2.3.2. Consumo de recursos de cómputo (memoria/GPU) \newline
   2.3.3. Proporción de consultas resueltas sin transmisión de datos a servicios de terceros \newline
   2.3.4. Superficies de exposición con control implementado, sobre el total identificado \\
\hline
\end{tabularx}

\vspace{0.2cm}
\noindent{\small \textit{Nota.} Los indicadores de la dimensión 2.2 se registran en los contextos $A_0$ y $A_1$ con el mismo instrumento, conforme al diseño preexperimental; los de las dimensiones 2.1 y 2.3 se obtienen sobre el banco de evaluación en las configuraciones de la Tabla \ref{tab:configuraciones}.}
\end{table}


\subsection{Métodos de obtención del conocimiento científico}

\subsubsection*{Métodos del nivel teórico}

\begin{itemize}
    \item \textbf{Método por índices.} Se utilizará este método para la construcción del marco teórico. A partir de un índice global con la base conceptual, elementos más importantes, tecnologías aplicadas, desafíos y limitaciones actuales relacionados con arquitecturas RAG y la autorreflexión, desglosados hasta lo más específico para sistematizar la literatura especializada.
    \item \textbf{Método de análisis documental.} Con una guía elaborada por el investigador, a partir de los indicadores de la operacionalización de las variables (Tabla \ref{tab:operacionalizacion}), se analizarán artículos científicos, documentación técnica de arquitecturas RAG y autorreflexión, y la documentación operativa interna de Dualbiz SRL relevante para la etapa de diagnóstico.
    \item \textbf{Método de modelación.} En la estructuración de la propuesta se aplica este método para encontrar las relaciones, cualidades y objetos que definen los componentes de recuperación, autorreflexión y generación, que integran la arquitectura y representarlos en el Modelo crítico híbrido, Critic-EmpRAG1.0.

    \item \textbf{Consulta a expertos (método Delphi).} Se aplica para la validación teórica del Modelo crítico híbrido, Critic-EmpRAG1.0, conforme al procedimiento de consulta iterativa y anónima a expertos formalizado por \textcite{DalkeyHelmer1963delphi} y sistematizado por \textcite{LinstoneTuroff1975delphi}. La selección de los participantes se realiza mediante el coeficiente de competencia, donde se expresa el nivel de conocimiento que el especialista declara sobre el tema y el grado de fundamentación de sus criterios; se integran a la consulta quienes alcanzan un coeficiente alto o medio alto. El grupo se conforma con expertos de dos perfiles complementarios: académicos e investigadores en Inteligencia Artificial y Procesamiento de Lenguaje Natural, y profesionales con experiencia en desarrollo de \textit{software} y soporte de TI. 
\end{itemize}

\subsubsection*{Métodos del nivel empírico}


\begin{itemize}
    \item \textbf{Método de observación.} Mediante una guía de observación elaborada por el autor, se aplicará para comparar el estado actual del acceso al Sistema de soporte de conocimiento empresarial interno en las áreas de TI y desarrollo en Dualbiz SRL. y posterior a la aplicación del modelo para verificar la transformación que se produce.
    
    \item \textbf{Método pre experimental.} Se aplicará para la validación empírica del Modelo crítico híbrido, Critic-EmpRAG1.0, en concordancia con la relación sujeto-objeto asumida por el paradigma. Se interviene deliberadamente introduciendo el modelo crítico híbrido en el acceso al Sistema de soporte de conocimiento empresarial interno de las áreas de TI y desarrollo de Dualbiz SRL., y se observa el efecto que dicha intervención produce (preprueba y posprueba con un solo grupo), contrastando el contexto inicial $A_0$, previo a la aplicación del modelo, con el contexto final $A_1$, posterior a su aplicación. Este proceso experimental, después de la sistematización a la literatura y la estructuración del modelo, se prevé que se realice de manera escalonada tomando en consideración la arquitectura RAG, y las variantes con la introducción del modelo crítico híbrido.
    
    Dada la temporalidad longitudinal prospectiva del estudio, las mediciones se registrarán hacia adelante en el tiempo, a lo largo de las iteraciones del proceso experimental, lo que permitirá observar la evolución del desempeño del prototipo y verificar la transformación en los indicadores propuestos en la operacionalización. Como componente complementario de esta misma validación empírica, y con el fin de establecer a qué elemento del modelo resulta atribuible el efecto observado, el desempeño del prototipo se contrasta además con el de configuraciones alternativas del sistema que difieren en un único componente cada vez, conforme se detalla en el epígrafe siguiente.

\end{itemize}

\subsubsection*{Niveles de la validación empírica}

La validación empírica comprende dos componentes que responden preguntas distintas y que, por tanto, operan sobre unidades de análisis e instrumentos distintos.

El \textbf{primer nivel, de carácter organizacional}, verifica la transformación del campo de acción mediante el preexperimento descrito: contrasta el contexto $A_0$, correspondiente al proceso manual de consulta vigente, con el contexto $A_1$, posterior al despliegue del prototipo. Su unidad de análisis es el personal de las áreas de TI y desarrollo, y sus indicadores son los de la variable dependiente 1, productividad del proceso de consulta, medidos con los mismos instrumentos en ambos momentos. Responde a la pregunta de si mejoró la productividad del proceso de consulta al conocimiento operativo en Dualbiz SRL.

El \textbf{segundo nivel, de carácter computacional}, verifica a qué mecanismo resulta atribuible dicha mejora. Su unidad de análisis es la consulta del banco de evaluación, y sus indicadores son los de la variable dependiente 2, acceso al conocimiento, en sus dimensiones de fundamentación y confiabilidad y de viabilidad bajo restricciones. Responde a la pregunta de si el efecto procede del crítico híbrido y no de la mera incorporación de un mecanismo de recuperación, del cambio de modelo o de la ejecución de ciclos adicionales de inferencia. Este segundo nivel replica el procedimiento comparativo que \textcite{jeong_improving_2024} emplearon para establecer el aporte de Self-BioRAG frente a Self-RAG, y resulta necesario porque el contraste $A_0 \rightarrow A_1$ modifica simultáneamente varios componentes y no permite, por sí solo, aislar el mecanismo que esta investigación postula.

Las configuraciones que integran el segundo nivel se presentan en la Tabla \ref{tab:configuraciones}. Todas ellas comparten el modelo base y el corpus documental; a partir de $C_2$ comparten además el mecanismo de recuperación y los parámetros de segmentación e indexación, de modo que cada contraste de interés difiere en un único componente.

\begin{table}[H]
\caption{Configuraciones del segundo nivel de la validación empírica}
\label{tab:configuraciones}
\setstretch{1}
\renewcommand{\arraystretch}{1.3}
\small
\begin{tabularx}{\textwidth}{|L{1.3cm}|L{4.6cm}|Y|}
\hline
\rowcolor{gray!25}
\textbf{Config.} & \textbf{Composición} & \textbf{Propósito} \\
\hline
$C_1$ & SLM sin recuperación & Línea base generativa: desempeño del modelo con su solo conocimiento paramétrico \\
\hline
$C_2$ & SLM + RAG convencional & Aporte de la recuperación sobre el corpus empresarial \\
\hline
$C_3$ & SLM + RAG + crítico paramétrico (Self-RAG) & Aporte de la autorreflexión cuya especificidad de dominio reside en los pesos (\cite{Asai2024_Selfrag}) \\
\hline
$C_4$ & SLM + RAG + Critic-EmpRAG1.0 & Sistema completo: crítico híbrido con juicio anclado en la evidencia recuperada \\
\hline
$C_{4a}$ & $C_4$ sin ciclo correctivo & Ablación del refinamiento iterativo \\
\hline
$C_{4b}$ & $C_4$ sin evidencia externa para el crítico & Ablación del componente no paramétrico del crítico \\
\hline
\end{tabularx}

\vspace{0.2cm}
\noindent{\small \textit{Nota.} Todas las configuraciones comparten el modelo base y el corpus documental; de $C_2$ en adelante comparten además el mecanismo de recuperación y los parámetros de segmentación e indexación. SLM: modelo de lenguaje pequeño; RAG: generación aumentada por recuperación.}
\end{table}


Los contrastes de interés son los siguientes. La diferencia entre $C_1$ y $C_2$ cuantifica el aporte de la recuperación; entre $C_2$ y $C_3$, el de la autorreflexión resuelta por vía paramétrica; y entre $C_3$ y $C_4$, el aporte específico del crítico híbrido, que constituye la novedad científica de esta investigación. Las configuraciones $C_{4a}$ y $C_{4b}$ operan como ablaciones del sistema completo: la primera suprime el ciclo correctivo y la segunda priva al crítico del acceso a la evidencia recuperada, conservando el resto de la arquitectura. Si $C_{4b}$ retrocede hacia el desempeño de $C_3$, queda establecido que la mejora procede del anclaje no paramétrico del juicio y no de la formulación de las instrucciones.

\subsection{Técnicas de investigación}

\begin{itemize}
    \item \textbf{Entrevista} al responsable de soporte y/o desarrollo de Dualbiz SRL. para determinar el estado actual del acceso al Sistema de soporte de conocimiento empresarial, recolectando los factores que se consideran relevantes para la mejora.
    \item \textbf{Encuesta} al personal de las áreas de TI y desarrollo de Dualbiz SRL, para obtener datos cuantificables sobre el tiempo y la efectividad del proceso actual de consulta al Sistema de soporte de conocimiento empresarial interno.
\end{itemize}

\subsection{Instrumentos de investigación}

\begin{itemize}
    \item Guía de análisis documental.
    \item Guía de entrevista al responsable de soporte y/o desarrollo de Dualbiz SRL.
    \item Cuestionario de encuesta al personal de soporte y/o desarrollo de Dualbiz SRL.
    \item Cuestionario de consulta a expertos.
    \item Guía de observación.
\end{itemize}

\subsection{Población y muestra}

La población y muestra coinciden al 100\%, corresponde al personal de las áreas de TI y desarrollo de Dualbiz SRL. incluyendo al responsable de dichas áreas, para un total de 18 personas.

\subsection{Tratamiento de la información}
Los datos cualitativos obtenidos de las entrevistas se analizarán mediante análisis de contenido, para identificar los factores críticos del proceso actual de acceso al conocimiento operativo y traducirlos en requisitos de diseño de la arquitectura que se propone. Los resultados de la consulta a expertos se procesan mediante el método Delphi. Los datos cuantitativos se procesan de manera diferenciada según el nivel de validación del que proceden.

En el primer nivel, los indicadores registrados en los contextos $A_0$ y $A_1$ provienen de las mismas dieciocho personas en dos momentos, por lo que constituyen mediciones pareadas. El contraste se realiza mediante la prueba $t$ de Student para muestras relacionadas, previa verificación del supuesto de normalidad de las diferencias mediante la prueba de Shapiro-Wilk, apropiada para muestras de tamaño reducido; de no satisfacerse dicho supuesto se aplica la prueba de Wilcoxon de rangos con signo. En ambos casos se reporta, junto al valor de probabilidad, el tamaño del efecto y su intervalo de confianza al 95\,\%, dado que un tamaño muestral de dieciocho unidades limita la potencia estadística disponible.

En el segundo nivel, cada configuración de la Tabla \ref{tab:configuraciones} se evalúa sobre el mismo banco de consultas, de modo que las observaciones también resultan pareadas por consulta. La comparación conjunta de las configuraciones se realiza mediante análisis de varianza de medidas repetidas cuando se satisfacen sus supuestos, o mediante la prueba de Friedman en caso contrario; los contrastes entre pares de interés se examinan posteriormente con corrección por comparaciones múltiples.

\subsection{Resultados Esperados}
Los resultados obtenidos sobre los datos del primer nivel contrastan la hipótesis 1 y los del segundo nivel, la hipótesis 2; cada una se acepta o se rechaza con independencia de la otra. Si se acepta la hipótesis 2 y se rechaza la 1, el mecanismo queda verificado sin que se demuestre su efecto sobre la productividad organizacional; si se acepta la 1 y se rechaza la 2, la mejora organizacional no resulta atribuible al crítico híbrido. Solo la aceptación de ambas sustenta la conclusión de que el Modelo crítico híbrido, Critic-EmpRAG1.0, mejora la productividad y el acceso al conocimiento por el mecanismo que esta investigación postula.


\subsection*{Aporte teórico}

En lo esencial, el Modelo crítico híbrido, Critic-EmpRAG1.0, extiende al componente crítico la distinción entre conocimiento paramétrico y no paramétrico formulada por \textcite{Lewis2020rag} para el componente generador. Separa conceptualmente dos atributos: \textbf{la competencia para juzgar}, de carácter genérico y paramétrico, heredada del modelo base; y el \textbf{conocimiento del dominio específico} (sobre el cual se juzga), de carácter no paramétrico y anclado en la evidencia recuperada en tiempo de inferencia. De esta separación se deriva la noción de crítico híbrido que articula la propuesta y que responde a cómo dotar de contexto de dominio específico al componente crítico sin reentrenarlo sobre datos del dominio propietario.

De esta separación se deriva un cuerpo de proposiciones que explican por qué y bajo qué condiciones el crítico híbrido mejora el juicio en dominio específico. Cada proposición se enuncia con su condición de aplicabilidad y con la observación que la refutaría dentro del diseño de validación (Tabla \ref{tab:configuraciones}).

\begin{itemize}
	\item \textbf{P1. Principio de separabilidad.} La competencia para juzgar y el conocimiento del dominio son atributos separables del componente crítico: el primero puede heredarse del modelo base sin ajuste, y el segundo puede suministrarse como evidencia en tiempo de inferencia. Se refuta si el crítico privado de la evidencia recuperada ($C_{4b}$) alcanza el desempeño del sistema completo ($C_4$), pues el conocimiento del dominio no estaría operando a través de la evidencia.
	\item \textbf{P2. Principio de sustitución.} Bajo restricciones que impiden el reentrenamiento, el anclaje del juicio en la evidencia recuperada sustituye al ajuste paramétrico como fuente de especificidad de dominio: un crítico híbrido alcanza en fundamentación y confiabilidad un desempeño igual o superior al de un crítico paramétrico de dominio general aplicado al mismo dominio. Se refuta si $C_3$ supera a $C_4$ en los indicadores de fundamentación y confiabilidad.
	\item \textbf{P3. Condición de cobertura.} El efecto enunciado en P2 está condicionado a que la evidencia recuperada contenga respaldo para la consulta. Cuando no lo contiene, el crítico híbrido carece de base para juzgar sobre el dominio y su comportamiento correcto es la abstención; por tanto, su ventaja debe manifestarse también como una mayor tasa de abstención correcta ante consultas sin respaldo. Se refuta si, en esas consultas, $C_4$ no se abstiene con mayor acierto que $C_2$ y $C_3$.
	\item \textbf{P4. Condición de capacidad.} La separabilidad enunciada en P1 exige una capacidad paramétrica mínima del modelo base: por debajo de ella, el crítico no puede evaluar la correspondencia entre la evidencia y la respuesta aunque disponga de la evidencia correcta. Esta condición se enuncia como línea de exploración de la tesis, ya que su contraste exige evaluar $C_3$ y $C_4$ con al menos dos modelos base de distinto tamaño dentro del rango de 1 a 13 mil millones de parámetros y queda sujeto a la infraestructura de cómputo disponible; si la ventaja de $C_4$ no varía con la capacidad del modelo, la condición se descarta.
\end{itemize}

En conjunto, las proposiciones establecen la relación causal que la hipótesis 2 somete a prueba: la disponibilidad de evidencia con respaldo (P3), procesada por un modelo base con capacidad suficiente (P4), permite que el crítico juzgue sobre el dominio sin conocerlo paramétricamente (P1), y ello produce respuestas más fundamentadas que las de un crítico ajustado por vía paramétrica (P2). El cuerpo integra así relaciones que hoy se enuncian de forma dispersa en la literatura ---que la autorreflexión pierde efectividad cuando el modelo carece de anclaje externo verificable (\cite{Shinn2023_Reflexion}), que la evaluación de la evidencia recuperada condiciona la calidad de la respuesta (\cite{Yan2024_CRAG}) y que el modelo no siempre reconoce los límites de su propio conocimiento (\cite{Jiang2024_KnowsWhat})--- y permite decidir, bajo restricciones dadas de \textit{hardware} y privacidad, qué mecanismo de autorreflexión resulta admisible en el dominio empresarial de la industria del \textit{software}.

En conjunto aportan un ordenamiento de los métodos, técnicas y despliegues del objeto de estudio en la intersección entre RAG, autorreflexión y modelos de lenguaje pequeños, con sus categorías, arquitecturas de referencia y vacíos de conocimiento identificados.

\subsection*{Aporte práctico}

El aporte práctico es el Modelo crítico híbrido, Critic-EmpRAG1.0, que se propone. Se concreta en los productos tangibles que transforman la situación descrita en los antecedentes y que quedan disponibles para la organización y para investigaciones posteriores:

\begin{itemize}
	\item \textbf{El prototipo funcional} que implementa el Modelo crítico híbrido, Critic-EmpRAG1.0, sobre modelos de lenguaje pequeños de código abierto, desplegado \textit{on-premise} en las áreas de TI y desarrollo de Dualbiz SRL., operando íntegramente dentro de las restricciones de privacidad de la información e infraestructura de cómputo de la empresa.
	\item \textbf{El banco de evaluación de conocimiento empresarial interno}, conformado por consultas representativas del proceso de soporte, sus respuestas esperadas y la evidencia documental de referencia que las sustenta, construido conforme a las políticas internas de confidencialidad de Dualbiz SRL.
	\item \textbf{El procedimiento replicable} de diseño y validación del componente crítico, que documenta las decisiones de diseño y su fundamento en la evidencia sistematizada, de modo que otras pequeñas y medianas empresas de desarrollo de \textit{software}, u otros dominios específicos con restricciones equivalentes, puedan reproducirlo sin repetir la fase exploratoria.
	\item \textbf{La transformación medible del campo de acción}, expresada en la mejora de la productividad y el acceso al Sistema de soporte de conocimiento empresarial interno, verificada mediante los indicadores de exactitud, fundamentación y eficiencia computacional.
\end{itemize}

\subsection*{Novedad científica}

La novedad científica de la investigación radica en dotar de especificidad de dominio al componente crítico de una arquitectura de autorreflexión sobre RAG en tiempo de inferencia y sin reentrenamiento sobre datos propietarios, mediante un crítico híbrido que ancla su juicio en la evidencia recuperada del corpus empresarial en lugar de hacerlo en parámetros ajustados con anotaciones de un modelo propietario. Esta formulación se distancia de las dos soluciones disponibles en la literatura precisamente en el punto que las hace inaplicables al contexto estudiado. Self-RAG (\cite{Asai2024_Selfrag}) opera sobre dominio general y deja abierta la adopción a dominios específicos; Self-BioRAG (\cite{jeong_improving_2024}) demuestra que esa adopción exige componentes específicos de dominio y la resuelve reentrenando el crítico con instrucciones anotadas mediante llamadas a un modelo propietario. 

Para una PyME tecnológica, ambas salidas resultan incompatibles de manera simultánea: la primera, porque la generalización deficiente conduce a recuperar documentos incorrectos y a emitir juicios inexactos sobre el dominio. La segunda, porque exponer el conocimiento operativo confidencial a un servicio de terceros, para su anotación, vulnera las restricciones de privacidad, y el reentrenamiento del crítico excede la infraestructura de cómputo disponible. La propuesta sortea ambas barreras al trasladar la especificidad de dominio del plano paramétrico al no paramétrico.

Según se evidencia en el SMS conducido por el investigador sobre 40 estudios primarios publicados entre 2020 y 2025 (\cite{Alcides2026}), no se identificaron estudios que satisfagan de forma concurrente las tres condiciones que definen este trabajo: la integración de autorreflexión y RAG sobre modelos de lenguaje pequeños en el rango de 1 a 13 mil millones de parámetros, su aplicación a un dominio empresarial privado real y su validación bajo restricciones efectivas de privacidad e infraestructura.

En consecuencia, resulta igualmente novedosa la naturaleza de la validación empírica que se propone: a diferencia de los antecedentes, que se verifican sobre conjuntos públicos de referencia y dentro de la distribución de entrenamiento del crítico, el Modelo crítico híbrido, Critic-EmpRAG1.0, se contrasta sobre conocimiento empresarial interno no publicado, ajeno a la distribución de entrenamiento de los modelos empleados, en un escenario operativo real de la industria del \textit{software}.

\subsection{Amenazas a la validez}

La validación empírica de una arquitectura RAG presenta riesgos documentados que conviene declarar y acotar antes de su ejecución. \textcite{Barnett2024_SevenFailurePoints} sistematizan, a partir de tres casos de estudio en dominios distintos, siete puntos de falla recurrentes en la ingeniería de estos sistemas, varios de los cuales inciden directamente sobre la validez de lo que aquí se propone medir.

\textbf{Validez interna.} El contraste entre los contextos $A_0$ y $A_1$ introduce de manera simultánea varios cambios: la existencia misma de un sistema automatizado, el modelo de lenguaje empleado, la estrategia de recuperación y el componente crítico. Los autores citados documentan además que el desempeño de un sistema RAG depende críticamente de la calibración del tamaño de fragmento, la estrategia de segmentación y el modelo de representación vectorial. De modo que, una mejora observada podría originarse en esas decisiones de configuración y no en el mecanismo que la investigación postula. Para que el efecto resulte atribuible al crítico híbrido, la evaluación del prototipo incorpora una condición de comparación construida sobre Self-RAG (\cite{Asai2024_Selfrag}), cuyo crítico resuelve la especificidad de dominio por vía paramétrica, manteniendo idénticos el modelo base, el corpus, el recuperador y los parámetros de segmentación. Se replica así el procedimiento comparativo que \textcite{jeong_improving_2024} emplearon para Self-BioRAG, de manera que la única diferencia entre ambas condiciones sea la fuente del conocimiento de dominio sobre el que el crítico emite su juicio.

\textbf{Validez de constructo.} El primer punto de falla identificado por \textcite{Barnett2024_SevenFailurePoints} advierte que un sistema RAG puede emitir una respuesta aparentemente válida ante preguntas que el corpus documental no está en condiciones de responder. Un banco de evaluación compuesto únicamente por consultas con respaldo documental mediría, en consecuencia, un escenario favorable y no verificaría la capacidad de abstención del modelo, que constituye uno de los referentes teóricos de esta propuesta. Por lo que, el banco de evaluación incorpora deliberadamente consultas sin respaldo en el corpus, y la abstención correcta ante ellas se registra como indicador. Se considera asimismo que el investigador interviene como diseñador del modelo, constructor del banco de evaluación y responsable de la medición; para atenuar el sesgo que dicha concurrencia introduce, la verdad de referencia del banco se establece con personal experto de la organización y se somete a revisión independiente sobre una submuestra.

\textbf{Validez de conclusión.} Los mismos autores concluyen que la validación de un sistema RAG solo resulta factible durante su operación y que su robustez evoluciona en lugar de quedar establecida en el diseño inicial. Dado que el alcance temporal de esta investigación excluye la operación productiva sostenida, los resultados se declaran válidos para la ventana de observación definida y no se extienden al comportamiento del sistema en régimen permanente, aspecto que se enuncia como línea de continuidad.

\textbf{Validez externa.} El estudio se realiza sobre un caso único, con una población de dieciocho personas que coincide con la muestra, lo que limita la potencia estadística y el alcance de la inferencia. Los resultados describen la transformación verificada en Dualbiz SRL; su transferencia a otras organizaciones con restricciones equivalentes de privacidad e infraestructura se plantea como hipótesis sujeta a validación posterior y no como resultado comprometido por el presente estudio.

\subsection{Procedimineto para la Investigación}
El procedimiento se organiza en dos momentos, que corresponden al diagnóstico y a la validación de la propuesta. En el diagnóstico, la información
obtenida con la guía de análisis documental, la guía de entrevista al responsable de las áreas de TI y desarrollo y la guía de observación se
procesa mediante análisis de contenido: se codifica según las dimensiones e indicadores de la Tabla \ref{tab:operacionalizacion} y se triangula
entre fuentes para identificar los factores críticos del proceso actual de consulta al conocimiento operativo en Dualbiz SRL, los cuales se traducen
en requisitos de diseño del Modelo crítico híbrido, Critic-EmpRAG1.0. Los datos del cuestionario de encuesta aplicado a las dieciocho personas se tabulan y
se resumen mediante estadística descriptiva por indicador (medidas de tendencia central y de dispersión) y constituyen, junto con los registros de
observación, la línea base del contexto inicial $A_0$. A partir de la misma documentación interna y con la participación del personal experto de la
organización se construye el banco de evaluación, que incluye consultas con y sin respaldo en el corpus documental.

En la validación, los criterios emitidos en el cuestionario de consulta a expertos se procesan por rondas sucesivas del método Delphi, en las que se
determina el grado de consenso sobre cada componente del modelo mediante el coeficiente de concordancia $W$ de Kendall y se incorporan los ajustes
antes de su implementación. Una vez desplegado el prototipo, se aplican nuevamente los mismos instrumentos para registrar el contexto final $A_1$, y
los pares de mediciones $A_0$--$A_1$ se procesan conforme al primer nivel descrito en el tratamiento de la información, con el fin de contrastar la
hipótesis 1. De manera complementaria, las respuestas que produce cada configuración de la Tabla \ref{tab:configuraciones} sobre el banco de
evaluación se califican con los indicadores de fundamentación y confiabilidad y de viabilidad bajo restricciones, y se procesan conforme al segundo
nivel para contrastar la hipótesis 2. Los resultados de ambos niveles se integran según los criterios de decisión enunciados en los resultados
esperados, y el análisis se realiza con herramientas de \textit{software} libre (Python o R), de modo que pueda reproducirse.


\subsection{Cronograma de investigación}\label{sec:cronograma}

La planificación tentativa de la investigación se expone mediante diagrama de Gantt en la Tabla \ref{tab:cronograma}, iniciando en octubre de 2026 y culminando en octubre de 2027.

\begin{table}[H]
\caption{Cronograma de la investigación (tentativo)}
\label{tab:cronograma}
\setstretch{1}
\renewcommand{\arraystretch}{1.3}
\setlength{\tabcolsep}{3pt}
\footnotesize
\begin{tabularx}{\textwidth}{|Y|*{13}{C{0.68cm}|}}
\hline
\rowcolor{gray!25}
 & \multicolumn{3}{c|}{\textbf{2026}} & \multicolumn{10}{c|}{\textbf{2027}} \\
\cline{2-14}
\rowcolor{gray!25}
\multirow{-2}{=}{\textbf{Actividades}} & \textbf{Oct} & \textbf{Nov} & \textbf{Dic} & \textbf{Ene} & \textbf{Feb} & \textbf{Mar} & \textbf{Abr} & \textbf{May} & \textbf{Jun} & \textbf{Jul} & \textbf{Ago} & \textbf{Sep} & \textbf{Oct} \\
\hline
Elaboración del perfil de tesis & X & & & & & & & & & & & & \\
\hline
Presentación del perfil de tesis & & X & & & & & & & & & & & \\
\hline
Revisión del perfil por el tribunal & & X & & & & & & & & & & & \\
\hline
Defensa del perfil ante el tribunal & & & X & & & & & & & & & & \\
\hline
Elaboración del marco teórico (profundizado) & & & & X & X & & & & & & & & \\
\hline
Diagnóstico del estado actual en Dualbiz SRL & & & & & & X & X & & & & & & \\
\hline
Diseño e implementación de la propuesta & & & & & & & & X & X & & & & \\
\hline
Validación de la propuesta & & & & & & & & & & X & X & & \\
\hline
Presentación del borrador del TFG & & & & & & & & & & & & X & \\
\hline
Revisión del borrador por el tribunal & & & & & & & & & & & & X & \\
\hline
Correcciones finales & & & & & & & & & & & & X & \\
\hline
Defensa final de tesis & & & & & & & & & & & & & X \\
\hline
\end{tabularx}

\vspace{0.2cm}
\noindent{\small \textit{Nota.} Cronograma tentativo, sujeto a ajuste según el avance real de la investigación.}
\end{table}


